\documentclass[11pt,a4paper]{article}

\usepackage[T1]{fontenc}
\usepackage{amsmath,amssymb,amsthm,mathtools}
\usepackage{booktabs}
\usepackage{enumitem}
\usepackage{geometry}
\usepackage{microtype}
\usepackage{listings}
\usepackage{hyperref}
\usepackage{array}

\geometry{margin=25mm}
\hypersetup{
  colorlinks=true,
  linkcolor=blue,
  citecolor=blue,
  urlcolor=blue,
  pdftitle={Cyclotomic-Correlation Reduction for Binary Circulants: The Exact Order-55 Maximum},
  pdfauthor={Qichao Wang and Daoyu Dong}
}
\urlstyle{same}
\setlength{\parindent}{0pt}
\setlength{\parskip}{5pt}
\setlength{\emergencystretch}{2em}
\lstset{
  basicstyle=\ttfamily\small,
  columns=fullflexible,
  keepspaces=true,
  breaklines=true,
  frame=single,
  framerule=0.3pt,
  xleftmargin=4pt,
  xrightmargin=4pt
}

\newtheorem{theorem}{Theorem}
\newtheorem{proposition}[theorem]{Proposition}
\newtheorem{lemma}[theorem]{Lemma}
\newtheorem{corollary}[theorem]{Corollary}
\theoremstyle{definition}
\newtheorem{definition}[theorem]{Definition}
\theoremstyle{remark}
\newtheorem*{remark}{Remark}
\DeclareMathOperator{\Res}{Res}
\DeclareMathOperator{\wt}{wt}
\newcommand{\Z}{\mathbb Z}
\newcommand{\D}{D_{01}(55)}
\newcommand{\M}{134694094094758395331307111329132}
\newcommand{\Pone}{2305843009213696591}
\newcommand{\Ptwo}{2305843009213697141}

\title{Cyclotomic--Correlation Reduction for Binary Circulants:\\The Exact Order-$55$ Maximum}
\author{
  Qichao Wang\\
  Hebei University of Technology
  \and
  Daoyu Dong\\
  University of Electronic Science and Technology of China
  \thanks{The authors contributed equally to this work.}
}
\date{}

\begin{document}
\maketitle

\begin{abstract}
Let $C(a)$ be the circulant matrix generated by a binary word
$a=(a_0,\ldots,a_{54})$.  We prove
\[
  D_{01}(55):=\max_{a\in\{0,1\}^{55}}\det C(a)
  =\M.
\]
Every maximizer has weight $28$, and all maximizers form one orbit under the
affine action $i\mapsto ui+t\pmod{55}$; the orbit has $2200$ elements and
trivial stabilizer.  The proof is a computer-assisted finite proof rather than
an unstructured search.  We first isolate cyclotomic and correlation identities
for binary circulants of odd order $pq$, where $p$ and $q$ are distinct primes.
For order $55=5\cdot11$, complement duality and exact spectral-moment bounds
reduce the problem to weights $25,26,27$ on the lower side of the complement
pairing.  Prime-factor folds give exact constraints on the order-$5$ and
order-$11$ cyclotomic sectors.  A complete recursion enumerates the remaining
integer correlation multisets, and a meet-in-the-middle lift enumerates every
binary word in each surviving margin-and-correlation fiber.  All proof-critical
comparisons use integer or rational arithmetic.  The final determinant is
verified by signed finite-field Fourier reconstruction, an exact integer matrix
determinant, and cyclotomic resultants.  The complete machine-readable proof
data and verification code are publicly available.
\end{abstract}

\section{Introduction}

The maximal determinant problem asks how large the determinant of a matrix can
be under a discrete restriction on its entries.  For circulant matrices the
Fourier transform diagonalizes the matrix~\cite{davis}, but the extremal problem remains
combinatorial because the generating word is discrete.  We write
\[
  D_{01}(n)=\max_{a\in\{0,1\}^n}|\det C(a)|.
\]
The general maximal determinant problem and its connections with Hadamard
matrices and design theory are surveyed by Browne, Egan, Hegarty and
\'O~Cath\'ain~\cite{browne-survey}.  Circulant binary determinants were studied
from the Hadamard-bound viewpoint by Maze and Parlier~\cite{maze-parlier}, and
closed-form evaluations for special binary circulant families were obtained by
Kravvaritis~\cite{kravvaritis} and Hariprasad~\cite{hariprasad}.  Brent and
Yedidia developed efficient exhaustive algorithms based on necklaces and
roots-of-unity determinant evaluation and determined the exact binary-circulant
maxima through order $53$ in the updated version of their work
\cite{brent-yedidia}.  A separate computation subsequently treated order $54$
\cite{aletheia54}.  To our knowledge, the exact value at order $55$ has not
previously been determined; the corresponding OEIS record also stops before
this case~\cite{oeis}.

Our purpose is not merely to append one integer to a table.  The factorization
$55=5\cdot11$ permits a useful reduction of the determinant problem to three
cyclotomic sectors, two low-dimensional folds, a finite correlation-profile
problem, and finally a binary fiber problem on the torus
$\Z/5\Z\times\Z/11\Z$.  The same sector identities hold for every odd
square-free semiprime order $pq$.  At order $55$ they are strong enough to make
the remaining finite search certifiable.

The main theorem is as follows.

\begin{theorem}[Exact order-$55$ maximum and classification]
\label{thm:main}
For binary circulant matrices of order $55$,
\[
  D_{01}(55)=\M.
\]
Every maximizing word has weight $28$.  The maximizing words form a single
orbit under
\[
  a_i\longmapsto a_{ui+t},\qquad
  u\in(\Z/55\Z)^\times,\quad t\in\Z/55\Z,
\]
and this orbit has $2200$ elements, hence trivial stabilizer.
\end{theorem}

The proof has four logically separate parts.  First, analytic identities reduce
all possible weights to three lower complementary weights.  Second, the
$5$- and $11$-folds restrict the possible cyclotomic sectors.  Third, an exact
recursion enumerates every integer correlation multiset that survives the
moment bounds, followed by all relevant ordered profiles modulo the unit
action.  Fourth, a meet-in-the-middle algorithm exhausts every binary fiber
corresponding to a retained target.  The central completeness statements are
proved in the text; the code executes these finite reductions and records their
outputs.

\section{Fourier and correlation identities at odd order}
\label{sec:fourier}

Let $n$ be odd, let $a=(a_0,\ldots,a_{n-1})\in\{0,1\}^n$, and put
\[
  k=\sum_{j=0}^{n-1}a_j,\qquad
  f_a(x)=\sum_{j=0}^{n-1}a_jx^j.
\]
Let $\zeta=\exp(2\pi i/n)$.  The eigenvalues of $C(a)$ are
$f_a(\zeta^r)$, $0\le r<n$.  For
\[
  q_r=|f_a(\zeta^r)|^2,\qquad 1\le r\le \frac{n-1}{2},
\]
complex conjugate pairing gives
\begin{equation}
\label{eq:det-fourier-general}
  \det C(a)=k\prod_{r=1}^{(n-1)/2}q_r\ge0.
\end{equation}
Thus the absolute value is redundant for actual binary words of odd order.

Define the periodic autocorrelations
\[
  c_s=\sum_{t=0}^{n-1}a_ta_{t+s},\qquad s\in\Z/n\Z.
\]
Then $c_0=k$ and $c_s=c_{n-s}$.  Counting ordered pairs of distinct support
elements gives
\begin{equation}
\label{eq:corrsum-general}
  2\sum_{s=1}^{(n-1)/2}c_s=k(k-1).
\end{equation}
Since $(c_s)$ is the inverse Fourier transform of the squared Fourier
magnitudes, Parseval gives
\begin{align}
  2\sum_{r=1}^{(n-1)/2}q_r&=nk-k^2,
  \label{eq:firstmoment-general}\\
  k^4+2\sum_{r=1}^{(n-1)/2}q_r^2
    &=n\left(k^2+2\sum_{s=1}^{(n-1)/2}c_s^2\right).
  \label{eq:secondmoment-general}
\end{align}
In particular, AM--GM yields
\begin{equation}
\label{eq:amgm-general}
  \det C(a)\le
  k\left(\frac{k(n-k)}{n-1}\right)^{(n-1)/2}.
\end{equation}

\subsection{Complement duality}

Let $\bar a=1-a$.  At every nontrivial $n$th root of unity,
$f_{\bar a}(\zeta^r)=-f_a(\zeta^r)$, while
$f_{\bar a}(1)=n-k$.  Since $n-1$ is even,
\begin{equation}
\label{eq:complement-general}
  k\det C(\bar a)=(n-k)\det C(a),\qquad 0<k<n.
\end{equation}
Therefore it suffices to analyze $1\le k\le(n-1)/2$; a lower-weight word
represents the complementary higher-weight determinant
\begin{equation}
\label{eq:compdet-general}
  \det C(\bar a)=(n-k)\prod_{r=1}^{(n-1)/2}q_r.
\end{equation}
Translations and multiplication of indices by a unit modulo $n$ preserve the
determinant.  These symmetries will be used only to choose representatives;
all orbit sizes are computed from actual stabilizers.

\section{Cyclotomic sectors for prime-product order}
\label{sec:pq}

We now isolate the structural feature used at order $55$.  Let
$n=pq$ with distinct odd primes $p$ and $q$.  For $m\mid n$ define the
$m$-fold vector
\[
  x^{(m)}_j=\sum_{\substack{0\le t<n\\t\equiv j\pmod m}}a_t,
  \qquad j\in\Z/m\Z,
\]
and its periodic autocorrelation
\[
  H^{(m)}_u=\sum_{j\in\Z/m\Z}x^{(m)}_j x^{(m)}_{j+u}.
\]
Directly from the definitions,
\begin{equation}
\label{eq:foldcorr-general}
  H^{(m)}_u=\sum_{\substack{s\in\Z/n\Z\\s\equiv u\pmod m}}c_s.
\end{equation}

For $m\in\{p,q,pq\}$ write
\[
 N_m(a)=\Res(f_a,\Phi_m)
       =\prod_{\substack{1\le r\le m\\(r,m)=1}}
          f_a(e^{2\pi ir/m}).
\]
Complex conjugation shows $N_m(a)\ge0$.  Grouping the roots of
$x^{pq}-1$ by exact order yields the following elementary factorization.

\begin{proposition}[Cyclotomic factorization]
\label{prop:cyclo-factor}
If $n=pq$ and $a$ has weight $k$, then
\[
  \det C(a)=kN_p(a)N_q(a)N_{pq}(a).
\]
\end{proposition}

\begin{proof}
The circulant determinant is $\prod_{\omega^n=1}f_a(\omega)$.  The root
$1$ contributes $f_a(1)=k$; all other roots have exact order $p$, $q$, or
$pq$, giving the three resultants.
\end{proof}

The folds also determine the first two power sums inside the prime-order
sectors.

\begin{proposition}[Prime-sector moment decomposition]
\label{prop:sector-moments}
For $m\in\{p,q\}$ let $\mathcal S_m$ be one representative from each
complex-conjugate pair of frequencies of exact order $m$, and define
\[
 A_m=\sum_{r\in\mathcal S_m}|f_a(\zeta^r)|^2,
 \qquad
 B_m=\sum_{r\in\mathcal S_m}|f_a(\zeta^r)|^4.
\]
Then
\begin{align}
 A_m&=\frac{m\sum_j(x^{(m)}_j)^2-k^2}{2},
 \label{eq:Am}\\
 B_m&=\frac{m\sum_u(H^{(m)}_u)^2-k^4}{2}.
 \label{eq:Bm}
\end{align}
If $A_{pq}$ and $B_{pq}$ denote the corresponding sums over primitive
$pq$-frequencies, then
\begin{align}
 A_{pq}
   &=\frac{pq\,k+k^2-p\sum_j(x^{(p)}_j)^2
                    -q\sum_j(x^{(q)}_j)^2}{2},
 \label{eq:Apq}\\
 B_{pq}
   &=\frac{pq\sum_{s\in\Z/pq\Z}c_s^2+k^4
       -p\sum_u(H^{(p)}_u)^2-q\sum_u(H^{(q)}_u)^2}{2}.
 \label{eq:Bpq}
\end{align}
\end{proposition}

\begin{proof}
Evaluation of $f_a$ at an $m$th root depends only on the folded polynomial
$\sum_j x^{(m)}_jz^j$.  Parseval on $\Z/m\Z$ gives
\[
 \sum_{r=0}^{m-1}\left|\sum_jx^{(m)}_j\zeta_m^{rj}\right|^2
   =m\sum_j(x^{(m)}_j)^2.
\]
Subtract the zero-frequency term $k^2$ and pair conjugates to obtain
\eqref{eq:Am}.  Applying Parseval to the folded autocorrelation gives
\[
 \sum_{r=0}^{m-1}\left|\sum_jx^{(m)}_j\zeta_m^{rj}\right|^4
   =m\sum_u(H^{(m)}_u)^2,
\]
which gives \eqref{eq:Bm}.  At order $pq$, all nonzero frequencies have
exact order $p$, $q$, or $pq$; subtracting the two prime sectors from
\eqref{eq:firstmoment-general} and \eqref{eq:secondmoment-general} yields
\eqref{eq:Apq} and \eqref{eq:Bpq}.
\end{proof}

A useful consequence is a necessary fold screen.  Put
$h_m=(m-1)/2$ and $h_{\rm rem}=(pq-1)/2-h_m$.  If a lower-weight word $a$
has complementary determinant at least $T$, then AM--GM on all sectors other
than $m$ gives
\begin{equation}
\label{eq:fold-necessary}
  (pq-k)N_m(a)
  \left(\frac{\frac{k(pq-k)}2-A_m}{h_{\rm rem}}\right)^{h_{\rm rem}}
  \ge T.
\end{equation}
This condition turns a determinant threshold into a finite constraint on a
short fold vector.

\section{Exact moment reduction at order \texorpdfstring{$55$}{55}}
\label{sec:moment55}

From now on $n=55$ and the target is
\[
  M=\M.
\]
This value is first supplied by an explicit binary witness, but the proof below
does not use any heuristic property of that witness beyond the exact integer
threshold $M$.

\subsection{Balanced-square and product bounds}

We use two small extremal facts repeatedly.  If nonnegative integers
$y_1,\ldots,y_h$ have sum $T=hb+r$ with $0\le r<h$, then
\begin{equation}
\label{eq:balanced-square}
  \sum_i y_i^2\ge \mathcal B(T,h):=(h-r)b^2+r(b+1)^2.
\end{equation}
Indeed, moving one unit from a larger entry to a smaller entry decreases the
square sum whenever their difference is at least two.

For the Fourier magnitudes we need an upper bound from their first and second
moments.

\begin{lemma}[Two-value product bound]
\label{lem:two-value}
Fix $h\ge2$, $S\ge0$ and $L\ge S^2/h$.  Among nonnegative real
$x_1,\ldots,x_h$ with $\sum_i x_i=S$ and $\sum_i x_i^2\ge L$, a positive
maximizer of $\prod_i x_i$ may be taken on the boundary
$\sum_i x_i^2=L$, and every positive stationary maximizer has at most two
distinct coordinate values.  If $x_i\le B$ is also imposed, every maximizer is
obtained by fixing some coordinates at $B$ and applying the same two-value
description to the remaining interior coordinates, together with the
zero-product boundary cases.
\end{lemma}

\begin{proof}
On the simplex $\sum x_i=S$, replacing two unequal positive coordinates by
values closer to their mean preserves the sum, decreases their square sum and
increases their product.  Thus a positive maximizer cannot have
$\sum x_i^2>L$ unless an upper boundary is active; after fixing active
coordinates the same argument applies recursively.  In the positive interior
of a fixed face, Lagrange multipliers give
\[
  \frac1{x_i}=\alpha+2\beta x_i,
\]
so every interior coordinate is a root of one quadratic equation.  Hence at
most two interior values occur.  There are only finitely many choices of the
number of upper-bound coordinates and of the multiplicity of either interior
root.
\end{proof}

When there is no cap and the smaller value occurs $m$ times, set
$\mu=S/h$ and $V=L-S^2/h$.  The two values are
\begin{equation}
\label{eq:two-values}
 x_- = \mu-\sqrt{\frac{V(h-m)}{hm}},\qquad
 x_+ = \mu+\sqrt{\frac{Vm}{h(h-m)}}.
\end{equation}
The implementation enumerates $m=1,\ldots,h-1$ and encloses every square root
by rational outward intervals at decimal scale $10^{50}$.  Thus no
floating-point comparison can decide a pruning step.

For a weight-$k$ word, let
\[
 S_1=\frac{k(55-k)}2,
 \qquad
 S_2(C)=\frac{55(k^2+2C)-k^4}{2},
 \qquad
 C=\sum_{s=1}^{27}c_s^2.
\]
Equations \eqref{eq:firstmoment-general}--\eqref{eq:secondmoment-general}
and Lemma~\ref{lem:two-value} give an exact upper bound
\[
  \det C(\bar a)
  \le(55-k)\,\mathcal P_{27}(S_1,S_2(C)),
\]
where $\mathcal P_h$ is the finite two-value maximum described above.  The
integer correlation sum is $k(k-1)/2$, so the minimum possible $C$ is
$\mathcal B(k(k-1)/2,27)$.

\begin{proposition}[Weight reduction]
\label{prop:weight-reduction}
If a binary circulant of order $55$ has determinant at least $M$, then after
complementing if necessary its lower weight belongs to $\{25,26,27\}$.
Moreover, a necessary condition is
\[
\begin{array}{c|c|c}
 k & \min\sum_{s=1}^{27}c_s^2 & \max\sum_{s=1}^{27}c_s^2\\
\hline
25&3336&3338\\
26&3913&3919\\
27&4563&4570.
\end{array}
\]
\end{proposition}

\begin{proof}
For each $1\le k\le27$ we evaluate exactly
\begin{equation}
\label{eq:Uk}
 U_k=\left\lfloor\min\left\{
 (55-k)\left(\frac{k(55-k)}{54}\right)^{27},
 (55-k)\mathcal P_{27}\bigl(S_1,S_2(C_{\min})\bigr)
 \right\}\right\rfloor,
\end{equation}
where $C_{\min}=\mathcal B(k(k-1)/2,27)$.  The complete values are listed in
Appendix~\ref{app:weights}; for $k\le24$ every $U_k<M$.  For
$k=25,26,27$, monotonicity of $S_2(C)$ and the same exact product bound give
the largest $C$ still capable of reaching $M$, namely $3338,3919,4570$.
The next values already have upper bounds
\[
\begin{array}{c|c}
 k & \text{first excluded upper bound}\\
\hline
25&133559927019518974061936346311463\\
26&129449029897850512940750031714722\\
27&134294603011751900149424725226911,
\end{array}
\]
all strictly below $M$.
\end{proof}

\subsection{A coefficient-rearrangement cap}
\label{subsec:cap}

The moment bound can be sharpened for a fixed correlation multiset.  Put
\[
  b=\left\lfloor\frac{k(k-1)}{54}\right\rfloor,
  \qquad d_s=c_s-b.
\]
For $1\le r\le27$,
\begin{equation}
\label{eq:q-deviation}
 q_r=k-b+\sum_{s=1}^{27}d_s\,2\cos\frac{2\pi rs}{55}.
\end{equation}
For a fixed multiset $\{d_s\}$, rearrangement bounds the right-hand side using
only the multiset of cosine coefficients.  At order $55$ there are three
coefficient multisets, represented by $r=1,5,11$, corresponding to
$\gcd(r,55)=1,5,11$.  Rational dyadic outward intervals for
$2\cos(2\pi t/55)$ are constructed from Machin's identity for $\pi$ and
Taylor remainders.  Pairing the sorted deviations with the appropriate
outward endpoints gives a certified common upper cap $B$ for every $q_r$.
Lemma~\ref{lem:two-value} with the box constraint $q_r\le B$ then eliminates
many correlation multisets without expanding their ordered permutations.

\section{Prime-factor fold reduction at order \texorpdfstring{$55$}{55}}
\label{sec:fold55}

For $m=5$ or $11$, a fold vector has length $m$, nonnegative integer entries,
sum $k$, and entrywise bound $55/m$.  Equation \eqref{eq:fold-necessary}
provides a necessary lower bound on its cyclotomic norm $N_m$.  This is a
finite composition problem.

\begin{lemma}[Completeness of the fold enumeration]
\label{lem:fold-complete}
For fixed $m\in\{5,11\}$ and active weight $k$, the fold generator visits a
representative of every affine orbit of fold vectors satisfying the entry
bounds and every square-sum threshold that can pass
\eqref{eq:fold-necessary}.  It retains exactly those representatives whose
cyclotomic norm meets the required threshold.
\end{lemma}

\begin{proof}
Every fold vector can be translated so that a minimum entry occurs in
coordinate $0$.  The recursion assigns the coordinates from left to right,
maintaining their partial sum and square sum.  The only prunings are necessary
conditions: the remaining entries cannot supply the missing total under the
entry cap, the balanced-square lower bound for the remaining sum already
exceeds the allowed square range, or the largest possible additional square
sum cannot reach the required range.  Hence a feasible vector has a surviving
prefix at every depth and is visited.  At a leaf the program compares the
vector against all $m(m-1)$ affine images $j\mapsto uj+t$ and retains the
lexicographically least representative.  Thus every orbit occurs exactly once.

The norm is evaluated at a primitive $m$th root in the prime field of size
$P_1=\Pone$.  Since $|\sum_jx_j\zeta_m^j|\le k$, we have the crude bound
$0\le N_m\le k^{m-1}<P_1$ for $k\le27$ and $m\le11$.  Therefore the finite
field residue uniquely equals the ordinary integer norm, so the threshold test
is exact.
\end{proof}

The retained representative counts are
\begin{center}
\begin{tabular}{@{}r|rr|rr@{}}
\toprule
&\multicolumn{2}{c|}{$m=5$}&\multicolumn{2}{c}{$m=11$}\\
$k$ & representatives & corr. signatures & representatives & corr. signatures\\
\midrule
25&17&19&321&1020\\
26&27&44&560&1765\\
27&26&45&660&2055\\
\bottomrule
\end{tabular}
\end{center}
Here a folded-correlation signature is the symmetric half of
$H^{(m)}$ from \eqref{eq:foldcorr-general}.  These signatures are the interface
between the fold stage and the full correlation stage.

\section{Complete enumeration of correlation profiles}
\label{sec:corrprofiles}

A lower-weight binary word determines $27$ integers
$c_1,\ldots,c_{27}$ with
\[
 \sum_{s=1}^{27}c_s=\frac{k(k-1)}2,
 \qquad
 \sum_{s=1}^{27}c_s^2\le C_{\max}(k).
\]
The first enumeration is over multisets, represented by nondecreasing tuples.
It is important that this stage has a mathematical coverage proof rather than
being described only by a program run.

\begin{lemma}[Exact multiset recursion]
\label{lem:multiset-recursion}
Fix integers $h,T,C$ and $0\le c_i\le k$.  The following recursion enumerates
exactly all nondecreasing $h$-tuples $(c_1,\ldots,c_h)$ satisfying
$\sum c_i=T$ and $\sum c_i^2\le C$:
\begin{enumerate}[leftmargin=2.4em,itemsep=2pt]
\item maintain a prefix, its last value $\ell$, remaining length $r$, remaining
sum $R$, and accumulated square sum $Q$;
\item prune if $R<\ell r$ or if
$Q+\mathcal B(R,r)>C$;
\item otherwise choose the next value
\[
  v=\ell,\ell+1,\ldots,
  \min\!\left(k,\left\lfloor R/r\right\rfloor\right)
\]
and recurse with $R-v$, $Q+v^2$, and last value $v$.
\end{enumerate}
\end{lemma}

\begin{proof}
Every emitted leaf is nondecreasing by construction, has the prescribed sum,
and passes the square bound.  Conversely, take any feasible nondecreasing
tuple and any one of its prefixes.  Its remaining entries are all at least the
last prefix value, so $R\ge\ell r$.  Their square sum is at least
$\mathcal B(R,r)$, so the second pruning condition also fails.  Its next value
is at least $\ell$, at most $k$, and at most $\lfloor R/r\rfloor$ because the
remaining tuple is nondecreasing.  Thus the recursion follows the tuple's
actual next value at every depth and reaches its leaf.  This proves both
soundness and completeness.
\end{proof}

For each emitted multiset, its number of ordered assignments is the exact
multinomial coefficient
\[
  \frac{27!}{\prod_v m_v!},
\]
where $m_v$ is the multiplicity of value $v$.  The capped spectral bound from
Section~\ref{subsec:cap} is then evaluated once per multiset.  If that bound is
below $M$, all of its ordered assignments are discarded at once.  Otherwise
every ordered assignment is generated explicitly, checked against both fold
signature sets, and canonicalized under the effective order-$20$ unit action
on the $27$ opposite-shift pairs.

This distinction explains the two different profile counts in Table
\ref{tab:profile-counts}.  The ``coarse ordered count'' is a combinatorial count
of every ordered profile represented by the complete multiset recursion.  The
``materialized assignments'' column counts only those ordered profiles whose
multiset survives the rigorous capped product bound; the difference is already
proved unable to reach $M$ and is not an unsearched remainder.

\begin{table}[ht]
\centering
\small
\begin{tabular}{@{}r|r r r r@{}}
\toprule
$k$ & coarse ordered count & materialized assignments & unit-canonical & lift targets\\
\midrule
25&407277&2925&55&55\\
26&32178654&816102&22375&281\\
27&6043753&6026203&155558&1980\\
\midrule
Total&38629684&6845230&177988&2316\\
\bottomrule
\end{tabular}
\caption{Exact correlation-profile counts.}
\label{tab:profile-counts}
\end{table}

For a formal ordered profile, define $q_r$ by
\eqref{eq:q-deviation}.  The values need not be nonnegative before binary
realizability is established.  Nevertheless
\[
  \Pi=(55-k)\prod_{r=1}^{27}q_r
\]
is an integer: Galois automorphisms permute the conjugate-pair factors, so the
product is a rational algebraic integer.  It is therefore legitimate to
compare the signed exact value of $\Pi$ with the incumbent.

The exact product is computed in the two prime fields
\[
  P_1=\Pone,\qquad P_2=\Ptwo,
\]
both congruent to $1\pmod{55}$, and reconstructed by signed CRT.  Uniqueness is
certified before reconstruction.  If
$S_2=\sum_{r=1}^{27}q_r^2$, Cauchy--AM--GM gives
\[
  |\Pi|\le(55-k)\left(\frac{S_2}{27}\right)^{27/2}.
\]
For every materialized multiset the verifier checks the integer inequality
\begin{equation}
\label{eq:crt-capacity}
  \bigl(2(55-k)\bigr)^2 S_2^{27}
     <(P_1P_2)^2 27^{27},
\end{equation}
which implies $2|\Pi|<P_1P_2$.  Hence the signed CRT value is unique.  The
remaining $2316$ targets are exactly those formal profiles whose exact product
is at least $M$ and whose two fold signatures are realizable by retained fold
vectors.

\begin{proposition}[Correlation-target coverage]
\label{prop:target-coverage}
Every lower-weight binary word of weight $k\in\{25,26,27\}$ whose complement
has determinant at least $M$ is affinely equivalent to a word whose full
correlation profile occurs among the $2316$ retained lift targets.
\end{proposition}

\begin{proof}
Proposition~\ref{prop:weight-reduction} places its correlation multiset inside
the input of Lemma~\ref{lem:multiset-recursion}.  Hence that multiset is
visited.  If its capped product bound were below $M$, its actual nonnegative
Fourier product would also be below $M$, a contradiction.  Therefore its
multiset is materialized, and the ordered correlation tuple induced by the
word is among the generated permutations.  Its $5$- and $11$-folds satisfy
\eqref{eq:fold-necessary}; Lemma~\ref{lem:fold-complete} therefore includes
their affine fold classes, so the tuple passes the fold-signature test after an
affine normalization.  Finally the unit-canonicalization selects one
representative from the corresponding orbit.  The exhaustive finite output of
this process contains $2316$ targets.
\end{proof}

\section{Exhaustive binary lifting on \texorpdfstring{$\Z/5\Z\times\Z/11\Z$}{Z/5Z x Z/11Z}}
\label{sec:lifts}

Use the explicit Chinese-remainder coordinate
\[
  (i,j)\in\Z/5\Z\times\Z/11\Z
  \longmapsto 11i+45j\pmod{55}.
\]
A binary word is uniquely represented by eleven column masks
$X_0,\ldots,X_{10}\in\{0,1\}^5$.  A lift target fixes a realizable $5$-fold,
a realizable $11$-fold, and a full symmetric correlation profile.  Enumerating
the compatible orientations of the two folds expands the $2316$ targets to
$16084$ finite lift tasks.

For one task, split the eleven columns into a left block of five and a right
block of six.  For a partial assignment $X$, define the seven-coordinate
signature
\[
  \sigma(X)=\bigl(R_0(X),\ldots,R_4(X),A_{11}(X),A_{22}(X)\bigr),
\]
where $R_i$ is the row count contributed by the partial columns and
$A_{11},A_{22}$ are the within-column correlation contributions for shifts
$11$ and $22$.  Every coordinate is additive and nonnegative.

The left recursion enumerates every five-bit mask of the prescribed column
weight and stores each surviving half by its exact packed signature.  The
right recursion does the same for the other six columns and requests the
componentwise complementary signature.  A partial branch is pruned only if a
row count or one of $A_{11},A_{22}$ already exceeds its final target.  After a
signature join, every other correlation $c_s$, $1\le s\le27$,
$s\ne11,22$, is checked exactly by a $55$-bit rotation and population count.

\begin{lemma}[Completeness of the meet-in-the-middle lift]
\label{lem:mitm}
For every lift task, the meet-in-the-middle procedure returns exactly the
binary words in the prescribed margin-and-correlation fiber.
\end{lemma}

\begin{proof}
Take any binary word in the fiber.  Its eleven column masks have a unique
left/right decomposition for the chosen $5+6$ split.  Along either recursion,
each partial row count and each partial within-column correlation is a sum of
nonnegative contributions and therefore cannot exceed the corresponding final
target of the complete word.  Thus neither half of this word is pruned.
Additivity gives
\[
  \sigma(X_{\rm L})+\sigma(X_{\rm R})=\sigma_{\rm target},
\]
so the exact signature lookup necessarily joins its two halves.  The joined
word then passes all remaining correlation checks because it belongs to the
fiber.  Conversely, a returned join has the prescribed row counts and shifts
$11,22$ by signature complementarity, has the prescribed column weights by
construction, and is returned only after every remaining independent
correlation has been checked.  Hence every returned word lies in the fiber.
\end{proof}

The result does not depend on performing a second implementation of the lift.
A complementary $6+5$ split is available as a consistency check, but
Lemma~\ref{lem:mitm} is the completeness argument.  Across all $16084$ tasks,
the recorded $5+6$ run checks $15487882832$ joined words and produces two
normalized lifted solutions; both lead to the same affine class at the global
maximum.

\section{Winner, uniqueness, and the near-difference-set structure}
\label{sec:winner}

A canonical maximizing word is
\[
\boxed{\texttt{0000000100011111011011001101011011010100011110101000111}},
\]
of weight $28$.  Its exact cyclotomic norms are
\[
  N_5=131,\qquad
  N_{11}=434039,\qquad
  N_{55}=84603917386870862636041,
\]
and
\[
  28N_5N_{11}N_{55}=\M.
\]
The same integer is obtained from the exact $55\times55$ determinant and from
signed CRT in the two fields $P_1,P_2$ above.

The winner has nonzero correlation distribution
\[
  13\ \text{four times},\qquad
  14\ \text{forty-six times},\qquad
  15\ \text{four times}.
\]
Its weight-$27$ complement therefore has correlations $12$ four times,
$13$ forty-six times, and $14$ four times.  Thus the optimal object is very
close to the perfectly flat correlation pattern of a cyclic
$(55,27,13)$ difference set.  Difference sets and their relation to periodic
autocorrelation are classical~\cite{baumert}; three-level autocorrelation and
almost difference sets are surveyed by Nowak~\cite{nowak}.

The exactly flat alternative is impossible for a short algebraic reason.
Suppose a cyclic $(55,27,13)$ difference set existed, and let
$\alpha=f_a(\zeta_5)$.  Flat correlations would give
$\alpha\overline\alpha=14$.  Since $\operatorname{ord}_5(2)=4$,
$\Phi_5$ is irreducible modulo $2$, so
$\Z[\zeta_5]/(2)$ is a field.  Reducing
$\alpha\overline\alpha=14$ modulo $2$ gives $\alpha=0$ in that field, hence
$\alpha=2\beta$ in $\Z[\zeta_5]$.  It follows that
$\beta\overline\beta=7/2$, impossible because a rational algebraic integer is
an integer.  This explains why the extremizer is forced away from the ideal
flat spectrum; related nonexistence questions for symmetric designs are part
of the design-theoretic background~\cite{braic}.

We can now finish the proof of Theorem~\ref{thm:main}.  Any determinant at
least $M$ has, by complement duality, a lower complementary weight
$1\le k\le27$.  Proposition~\ref{prop:weight-reduction} excludes $k\le24$.
For $k=25,26,27$, Proposition~\ref{prop:target-coverage} moves the word into
one of the $2316$ target orbits, and Lemma~\ref{lem:mitm} guarantees that its
binary representative occurs in the exhaustive lift output.  The two
normalized output solutions have maximum complementary determinant $M$ and
the same affine canonical representative.  Direct enumeration of the affine
action on the displayed weight-$28$ word produces $2200$ distinct words.
Since $55\varphi(55)=2200$, the stabilizer is trivial, and there is exactly one
maximizing affine class.

\section{Reproducibility and proof artifacts}
\label{sec:repro}

The code and machine-readable proof data are available at
\begin{center}
\url{https://github.com/dongxuelian2/order55-binary-circulant-maxdet}.
\end{center}
The repository is a reproducibility companion; the mathematical coverage
arguments are Lemmas~\ref{lem:fold-complete}, \ref{lem:multiset-recursion},
and \ref{lem:mitm} together with Proposition~\ref{prop:target-coverage}.
The stored certificate has status \texttt{COMPLETE}, records the $2316$
correlation targets, $16084$ lift tasks and $15487882832$ joined words, and
identifies every component by SHA-256.  The SHA-256 digest of the certificate
manifest recorded by \texttt{global.json} is
\begin{center}
\small\ttfamily
b73a12c5b8502ac5cba20f5164ff351ba5df1a208d7cd2e67f45c359d7eedfc5.
\end{center}

The proof obligations and their executable counterparts are summarized in
Table~\ref{tab:proof-map}.  Historical host-specific absolute paths stored in
\texttt{build.json} are nonsemantic; source hashes, compiler flags, generated
files, exact counts, and mathematical outputs are the reproducibility data.

\begin{table}[ht]
\centering
\small
\begin{tabular}{@{}>{\raggedright\arraybackslash}p{0.29\textwidth}
                  >{\raggedright\arraybackslash}p{0.34\textwidth}
                  >{\raggedright\arraybackslash}p{0.28\textwidth}@{}}
\toprule
Proof obligation & Executable/source & Stored artifact\\
\midrule
weight and moment exclusion
& \texttt{src/maxdet/order55.py}; verifier
& \texttt{weight\_screen.json}\\
fold-orbit and norm enumeration
& \texttt{native/order55\_profiles.cpp}
& \texttt{profiles\_m*\_k*.jsonl}, \texttt{foldcorr\_*}\\
correlation multiset coverage and exact products
& native correlation-screen source; verifier recursion
& correlation-profile JSON; target files\\
binary fiber exhaustion
& native torus-lift source
& lift targets, coverage JSON, word files\\
final determinant and affine class
& verifier; exact matrix determinant; resultants
& \texttt{winner.json}, \texttt{global.json}\\
\bottomrule
\end{tabular}
\caption{Mapping from mathematical proof obligations to executable checks and
immutable proof artifacts.}
\label{tab:proof-map}
\end{table}

From the repository root, the stored proof data can be checked with
\begin{lstlisting}
python -m pip install -e '.[dev]'
python scripts/verify_order55_global.py
python -m pytest
\end{lstlisting}
The full finite enumeration can be regenerated with
\begin{lstlisting}
python scripts/build_order55_global.py --workers 8
python scripts/verify_order55_global.py
\end{lstlisting}
A C++20 compiler with unsigned \texttt{\_\_int128} support is required for the
native programs.  The optional \texttt{--full-lifts} mode repeats the lift with
the complementary $6+5$ split as a consistency check; it is not an additional
assumption in the proof.

\paragraph{Use of generative AI.}
Generative-AI tools were used during exploratory derivation, software
development, literature organization, and manuscript editing.  The authors
are responsible for all mathematical statements, code, and reported
computational results.  The final claims are supported by the arguments in
this manuscript and by the exact reproducibility artifacts described above.

\section{Discussion}

The order-$55$ computation illustrates a broader strategy for composite odd
circulant orders.  A direct necklace search treats the word as a monolithic
object.  The present approach first exposes the algebraic decomposition of the
Fourier spectrum, transfers the small cyclotomic sectors to short integer fold
vectors, uses moment information to constrain the full autocorrelation, and
postpones binary realizability until only a small set of fibers remains.  The
prime-product identities in Section~\ref{sec:pq} therefore separate the
order-specific numerical work from the reusable mathematics.

At order $55$, the winning correlation distribution is also informative: the
maximizer lies only a small defect away from the flat pattern associated with
a cyclic difference set, while the flat pattern itself is algebraically
forbidden.  Understanding whether this ``near design forced by a norm
obstruction'' phenomenon persists at other semiprime orders is a natural
structural question.  A second direction is algorithmic: one may ask for
conditions on $p,q$ under which the fold inequalities alone force a bounded
number of correlation multisets as $pq$ grows.  These questions are separate
from the finite verification of Theorem~\ref{thm:main}.

\appendix
\section{Complete weight screen}
\label{app:weights}

For transparency, Table~\ref{tab:allweights} lists the exact integer upper
bounds $U_k$ from \eqref{eq:Uk} for every lower complementary weight.  The
threshold is $M=134694094094758395331307111329132$.

\begin{table}[ht]
\centering
\scriptsize
\begin{tabular}{@{}r r@{\qquad}r r@{}}
\toprule
$k$ & $U_k$ & $k$ & $U_k$\\
\midrule
1&54&15&477591992253406544595615911884\\
2&395500979&16&947259933198935209412301505977\\
3&8310295299842&17&4192213634834240365480685545746\\
4&10869301568561838&18&5886783493739880805199443842870\\
5&3198970462262506504&19&12060322226419659758794582060128\\
6&399271202794134987472&20&34343188355470553561485146833558\\
7&32417513500367909529143&21&41532355024762070616242898837945\\
8&1903907609822691882312392&22&55595805954299054538842074345420\\
9&10258882626790513846158736&23&80909340172778266526436100785913\\
10&110320100452342521380555758&24&118322613332906266481408544824787\\
11&1831007961001001960240304341&25&165874915633519586770334322051012\\
12&4836960042417220465399169805&26&214709407622211262053744593911697\\
13&36982789534293066100191485862&27&246953391382495871652459563712512\\
14&97977612253598905895892393515&&\\
\bottomrule
\end{tabular}
\caption{Exact global upper bounds for all lower complementary weights.  The
first weight not excluded by the target $M$ is $k=25$.}
\label{tab:allweights}
\end{table}

\section{Finite enumeration specification}
\label{app:spec}

For clarity, we record the finite objects that a reimplementation must cover.
The correlation stage starts from all nondecreasing $27$-tuples of integers in
$[0,k]$ satisfying the exact sum in \eqref{eq:corrsum-general} and the
weight-specific square bound in Proposition~\ref{prop:weight-reduction}.
Lemma~\ref{lem:multiset-recursion} is a complete recursive specification of
that set.  For every surviving multiset, all distinct permutations are
visited.  A permutation can be discarded only by one of three conditions:
(1) its $5$-fold correlation signature is absent from the complete fold set;
(2) its $11$-fold signature is absent; or (3) its exact formal complementary
product is below $M$.  Unit canonicalization removes duplicates but does not
remove an orbit.

For each retained target, every compatible orientation of the two prime fold
vectors is included in the task list.  A task contains the five row margins,
the eleven column margins, and all $55$ correlations.  Lemma~\ref{lem:mitm}
then completely specifies the binary fiber computation.  The verifier checks
that the target list expands to exactly $16084$ distinct task descriptors and
that every descriptor has one stored lift record.  In this sense the
certificate is a finite transcript of the reductions proved in the paper,
not a substitute for their completeness proofs.

\begin{thebibliography}{99}

\bibitem{browne-survey}
P.~Browne, R.~Egan, F.~Hegarty, and P.~\'O~Cath\'ain,
\newblock A survey of the Hadamard maximal determinant problem,
\newblock \emph{Electronic Journal of Combinatorics} \textbf{28}(4) (2021),
P4.41.
\newblock \url{https://doi.org/10.37236/10367}

\bibitem{davis}
P.~J. Davis,
\newblock \emph{Circulant Matrices}, 2nd ed.,
\newblock AMS Chelsea Publishing, 1994.

\bibitem{maze-parlier}
G.~Maze and H.~Parlier,
\newblock Determinants of binary circulant matrices,
\newblock in \emph{Proceedings of the 2004 IEEE International Symposium on
Information Theory}, 2004, p.~120.
\newblock \url{https://doi.org/10.1109/ISIT.2004.1365161}

\bibitem{kravvaritis}
C.~Kravvaritis,
\newblock Determinant evaluations for binary circulant matrices,
\newblock \emph{Special Matrices} \textbf{2} (2014), 187--199.
\newblock \url{https://doi.org/10.2478/spma-2014-0019}

\bibitem{hariprasad}
M.~Hariprasad,
\newblock Determinant of binary circulant matrices,
\newblock \emph{Special Matrices} \textbf{7} (2019), 92--94.
\newblock \url{https://doi.org/10.1515/spma-2019-0008}

\bibitem{brent-yedidia}
R.~P. Brent and A.~B. Yedidia,
\newblock Computation of maximal determinants of binary circulant matrices,
\newblock \emph{Journal of Integer Sequences} \textbf{21} (2018), Article
18.5.6; updated version arXiv:1801.00399.
\newblock \url{https://arxiv.org/abs/1801.00399}

\bibitem{baumert}
L.~D. Baumert,
\newblock \emph{Cyclic Difference Sets},
\newblock Lecture Notes in Mathematics, vol.~182, Springer, 1971.
\newblock \url{https://doi.org/10.1007/BFb0061260}

\bibitem{nowak}
K.~Nowak,
\newblock A survey on almost difference sets,
\newblock arXiv:1409.0114, 2014.
\newblock \url{https://arxiv.org/abs/1409.0114}

\bibitem{braic}
S.~Brai{\v c},
\newblock Primitive symmetric designs with at most 255 points,
\newblock \emph{Glasnik Matemati\v{c}ki} \textbf{45}(2) (2010), 291--305.
\newblock \url{https://doi.org/10.3336/gm.45.2.01}

\bibitem{oeis}
OEIS Foundation Inc.,
\newblock A086432: Maximum determinant of binary circulants,
\newblock OEIS entry and b-file.
\newblock \url{https://oeis.org/A086432}

\bibitem{aletheia54}
K.~Terauchi,
\newblock \emph{ALETHEIA-54}, 2026.
\newblock \url{https://github.com/ruu413/aletheia-54}

\end{thebibliography}

\end{document}
